\section{Discussion}\label{sec:discussion}
The two negative timing outcomes are consistent with the distinction between
algebraic size and numerical work. An original-row formulation can save
auxiliaries while adding a more expensive proximal operation, a threshold
sort, a source conversion or an acceptance calculation. Sparse factorization
also depends on fill, ordering, numeric scaling and reuse. Variable counts
alone do not measure these effects. The small C01 formulation comparison
therefore explains representation but cannot establish why a complete
repeated-forecast operation is faster or slower.

The strongest practical comparison is an equivalent objective with a capable
reference used according to its documented interface. Directional row
splitting is especially important because it reuses a published grouped
projection method to express unequal directions. Compact epigraphs offer
another exact route. Removing those references would test a weaker question.
Likewise, allowing only the candidate to preserve its workspace would build
an artificial advantage into repeated execution. Online active-set methods
have long exploited parametric structure~\cite{qpOASES}; a reuse mechanism
is not in itself evidence of a new algorithmic contribution.

The grouped threshold calculation has clear mathematical value even when
its runtime prediction fails. It exposes the actual directional metric,
zero-price behavior and joint dual budget. Those details make a correctness
check possible and clarify why a dense metric or overlapping groups require
additional analysis. Existing weighted projection algorithms and simplex/
norm-ball work provide relevant algorithmic background~\cite{weighted,duchi,condat}.
The implemented threshold is not presented as the first use of sorting,
orthant behavior or scalar active-prefix equations.

Numerical policy and physical policy must also remain distinct. A stronger
temperature price can reduce a comfort departure by changing the optimized
tradeoff, not by improving numerical accuracy. An oracle-driven predictive
trajectory can isolate solver agreement on common inputs, not demonstrate
that separately optimized controllers have equal field behavior. For energy
examples, comparing a mixed group/Huber optimum with an authored separable
variant describes a policy change. These distinctions limit what a
practitioner should infer from cross-case objective tables and interactive
cost edits.

The results suggest choosing an admitted exact formulation based on the
actual problem regime and measured complete-operation requirements, rather
than adopting the smallest principal matrix by default. Where hard policy
is immutable, a method that cannot preserve it should be refused. Where
uncertainty about physical coefficients dominates numerical error, tighter
solver tolerances do not certify a better operational model. A future
computational prediction should specify a new mechanism and evaluation
prospectively; the existing held-out data cannot be retuned into evidence
for a revised threshold, population or favorable subset.

The system's inspectability remains useful despite the failed performance
criteria. Original-space quantities, exact compatibility and retained
unsuccessful outcomes let a reader distinguish an implementation defect
from an objective mismatch or a negative computational finding. They do
not establish adoption, realized savings or safety. Those questions require
different data, deployment conditions and evidence than the synthetic
optimization studies described here.
